% SGA 3, Expose V, Loop 1 English reconstruction.
% Scope: Section 5(b), general-case construction through the first paragraph
% on the effective equivalence relation R; ends with local p. 18.
% Authority: Polo--Gille Exp5-13oct24.pdf, component-local p. 18
% (combined-reader p. 280), re-edited page header 268; internal marker 271.
% Jacob Reinhold's CC BY 4.0 Markdown is comparison/drafting material only.

\phantomsection\label{proof:sga3-V-5-b-general}%
由\hyperref[par:sga3-V-5-a]{步骤~a)}, 可假设 \(d_1\) 局部自由且秩为 \(n\).
设 \((Y,p)\) 是所有赋环空间范畴中 \((d_0,d_1)\) 的余核.
\hyperref[proof:sga3-V-4.1-c]{4(c)} 末尾的论证表明,
为证明\hyperref[thm:sga3-V-4.1-i]{定理~4.1(i)},
只须证明 \(Y\) 是概形且 \(p:X_0\to Y\) 是概形态射.
由引理~\hyperref[sga3:exp5:lemma:1.2]{1.2}, 此问题在 \(Y\) 上是局部的.
设 \(y\in Y\), 并选取 \(x\in X_0\), 使 \(p(x)=y\).
若 \(x\) 有一个饱和仿射开邻域 \(U\), 则\hyperref[sec:sga3-V-4]{第~4 节}
表明 \(p(U)\) 是 \(Y\) 的仿射开集,
且 \(p|_U:U\to p(U)\) 是概形态射.
因此只须证明每个 \(x\in X_0\) 都有饱和仿射开邻域.
如下进行; 论证取自 SGA~1, VIII, 推论~7.6.

\SGAThreeDiagram[0.90\textwidth]
  {figures/exp5/p268_section5_affine_saturated_chain.png}
  {用于构造 \(x\) 的饱和仿射邻域的嵌套开集.}
  {fig:sga3-exp5-p268-affine-saturated-chain}
  {Polo--Gille Exp5-13oct24.pdf, 分篇第 18 页;
  合订本第 280 页; 重版本页眉页码 268}

由\hyperref[thm:sga3-V-4.1-b]{条件 4.1(b)},
有包含 \(d_1(d_0^{-1}(x))\) 的仿射开集 \(V\subseteq X_0\).%
\footnote{\label{note:sga3-V-28}编者注:
\(d_1(d_0^{-1}(x))=d_0(d_1^{-1}(x))\);
见\hyperref[thm:sga3-V-4.1]{定理~4.1} 中的\hyperref[note:sga3-V-16]{编者注~16}.}
置 \(F=X_0\setminus V\). 由于 \(d_1\) 整,
\(d_1(d_0^{-1}(F))\) 是闭集, 且
\[
  V'=X_0\setminus d_1(d_0^{-1}(F))
\]
是包含在 \(V\) 中的最大饱和开集.
集合 \(V'\) 是有限集 \(d_1(d_0^{-1}(x))\) 的邻域.
因此存在 \(V\) 的结构层的一个截面 \(f\), 它在
\(V\setminus V'\) 上为零, 且 \(d_1(d_0^{-1}(x))\)
包含在主开集 \(V_f\subseteq V\) 中.
我们将证明包含在 \(V_f\) 中的最大饱和开集 \((V_f)'\) 仿射;
它就会是所求的邻域.

置 \(Z(f)=V'\setminus V_f\). 则 \(d_0^{-1}(Z(f))\)
是
\[
  d_0^{-1}(V')=d_1^{-1}(V')
\]
中使 \(d_0^\ast(f)\) 为零的点所成的集合.
\(d_1\) 从这一共同逆像到 \(V'\) 诱导的态射局部自由且秩为 \(n\).%
\footnote{\label{note:sga3-V-29}编者注: 对
引理~\hyperref[lem:sga3-V-4.1.1]{4.1.1} 的引用是新增的;
见 \hyperref[bib:sga3-V-DG70]{[DG70]}, III, \S5.2, 第~313 页.}
因此引理~\hyperref[lem:sga3-V-4.1.1]{4.1.1} 表明,
\(d_1(d_0^{-1}(Z(f)))\) 恰好是 \(d_0^\ast(f)\)
相对于 \(d_1\) 的范数 \(N\) 的零点集. 可知
\[
  (V_f)'=V'\setminus d_1(d_0^{-1}(Z(f)))
\]
是 \(V_f\) 中使 \(N\) 不为零的点所成的集合.
因此 \((V_f)'\) 仿射.

这就证明了\hyperref[thm:sga3-V-4.1-i]{定理~4.1(i)};
断言 \hyperref[thm:sga3-V-4.1-ii]{(ii)},
\hyperref[thm:sga3-V-4.1-iii]{(iii)} 及
\hyperref[thm:sga3-V-4.1-iv]{(iv)} 的第一部分随即显然.
还须证明 \hyperref[thm:sga3-V-4.1-iv]{(iv)} 末尾陈述的推论;
见 \hyperref[bib:sga3-V-Ray67a]{[Ray67a]}, 定理~1(iii).

由假设, 群胚 \(X_\ast\) 来自一个等价关系
\[
  i:R\longrightarrow X_0\times X_0,
\]
其中 \(i\) 是浸入 (见\hyperref[note:sga3-V-19]{编者注~19}).
已证明 \(R\) 有效 (见 \SGARefLink{sga3:iv:definition:3:3:2}{第 IV 篇, 3.3.2}), 且
\[
  p:X_0\longrightarrow Y=X_0/R
\]
满射且有限局部自由; 特别地, 它忠实平坦且有限呈示.

% Translation coverage: entire assigned component; adjoining components retain the original statement/proof seams.
% Source lines: 85; translated body lines before this audit comment: 82.
% Math preserved: 49 inline formulas, 5 display blocks, 0 equation environments; mathematical prose translated where present.
% Labels: 3/3; hyperrefs: 18/18; SGARefLinks: 1/1; inline targets: 0/0; footnotes: 2/2.
% Figure arguments preserved; all rendered captions and natural-language source locators translated.
% Unresolved TODOs: 0; compilation and visual inspection remain with the volume integrator.
